We train WavLM-based \cite{chen2022wavlm} encoder-only models using synthetic LibriSpeech mixtures from FastMSS~\cite{polokfastmss} and scraped multi-talker YODAS data~\cite{li2023yodas}. NSF1 and AMI SDM dev and eval sets are out-of-domain to reflect the intended deployment scenario.

FastMSS contains 150,000 two-minute, single-channel, 16-kHz mixtures with 2--8 speakers, at most three concurrent speakers, and optional reverberation and noise. Mixtures are divided into utterance groups at silences of at least 1\,s, excluding groups longer than 80\,s. YODAS contributes 5,360 weakly labeled recordings selected as likely English multi-talker speech; see~\cite{polok2027other}. Segments are normalized to LibriSpeech-style uppercase text, grouped across pauses of at most 1\,s, and divided into 60\,s groups with 2--8 speakers. FastMSS and YODAS are sampled equally and share a 5,000-unit unigram BPE model trained on normalized LibriSpeech-960 and YODAS transcripts. We use 5k BPE units because our focus is long-form word-level alignment. We train an 8-speaker WavLM speaker-attributed shuffle model and a token-only model on 4 V100 GPUs for 72k iterations, using 4 gradient-accumulation steps for an effective batch size of 2880\,s. Other settings follow~\cite{wiesner2026modeling}. For long-form decoding, audio is processed in overlapping 30\,s chunks where the center 20\,s of the output is passed to the alignment graph. Speaker labels are resolved across chunks as in Section \ref{sec:longform_tricks}.

We compare against: (i) an audio-free baseline that derives word times from coarse segment boundaries assuming a uniform token rate; (ii) Montreal Forced Aligner using \texttt{english\_mfa} v3.1.0 with wider beams (\texttt{beam}$=40$, \texttt{retry\_beam}$=160$) to recover recordings that fail with default settings; and (iii) the Wav2Vec2 CTC aligner used for post-hoc alignment in WhisperX~\cite{bain2023whisperx}. We also compare with the model from~\cite{wiesner2026modeling} on its more challenging \texttt{dev3spk} setting, which provides a synthetic in-domain evaluation.

We use the three alignment metrics from~\cite{wiesner2026modeling}. Intersection over Union (IoU)~\cite{jaccard1901etude} and mean per-recording Boundary Error (BE)~\cite{huang2024less} measure boundary quality. Normalized Kendall-$\tau$ distance ($\tau$) measures the average per-token number of adjacent swaps needed to recover the reference token order. All alignments are calibrated for start and duration bias, with offsets learned on dev applied to eval. For \texttt{dev3spk}, we calibrate on the first half and report results on the second, following~\cite{wiesner2026modeling}. Speaker attribution is evaluated using word diarization error rate (WDER) \cite{shafey19_interspeech}, computed exactly over all tokens since the reference transcript is known.
 
%We report 0-collar diarization error rate (DER) only when re-purposing the multi-talker ASR system for diarization.


